\section{Analysis}

\subsection{Why Does Zero-Shot SVD Work?}

The pretrained weight $W_Q$ has an intrinsic low-rank structure: for a 
768$\times$768 matrix trained on natural language, the singular values decay 
rapidly. \textbf{Figure 8.1} shows the singular value spectrum for a middle 
DistilBERT layer:

\begin{itemize}
\item Top 128 singular values capture \textbf{{{TOP128\_ENERGY}}\%} of energy.
\item Top 192 capture \textbf{{{TOP192\_ENERGY}}\%}.

\end{itemize}
This means $k = 128$ gives an excellent approximation without any 
calibration.

\subsection{Per-Layer Analysis}

\textbf{Table 8.1: Reconstruction error by layer.}

\begin{table}[H]
\centering
\small
\begin{tabular}{llll}
\toprule
\textbf{Layer} & \textbf{Q err} & \textbf{K err} & \textbf{V err} \\
\midrule
0 & {{D\_L0\_Q}} & {{D\_L0\_K}} & {{D\_L0\_V}} \\
3 & {{D\_L3\_Q}} & {{D\_L3\_K}} & {{D\_L3\_V}} \\
5 & {{D\_L5\_Q}} & {{D\_L5\_K}} & {{D\_L5\_V}} \\
\bottomrule
\end{tabular}
\end{table}

\textbf{Observation.} Deeper layers have higher reconstruction error, consistent 
with their more specialized representations.

\subsection{Where Does Accuracy Loss Come From?}

\subsubsection{Reconstruction Error}

The factorized Q/K/V matrices introduce approximation error. Our Frobenius 
error is \textbf{{{D\_REC\_AVG}}\%} on average.

\subsubsection{Interaction with $W_O$}

$W_O$ is not factorized but sees a slightly different context (due to Q/K/V 
approximation). This is the dominant source of accuracy loss.

\subsubsection{Attention Sink Effects}

For decoders, the attention sink at position 0 becomes slightly weaker. This 
may contribute to the GPT-2 drop of {{G\_DELTA}}\%.

\subsection{Comparison to Vanilla SVD}

\textbf{Table 8.2: SOA Fused vs vanilla SVD.}

\begin{table}[H]
\centering
\small
\begin{tabular}{lll}
\toprule
\textbf{Model} & \textbf{Vanilla SVD (no LoRA)} & \textbf{SOA Fused (with LoRA)} \\
\midrule
DistilBERT & {{D\_VANILLA}} & \textbf{{{D\_ACC}}} \\
RoBERTa & {{R\_VANILLA}} & \textbf{{{R\_ACC}}} \\
GPT-2 & {{G\_VANILLA}} & \textbf{{{G\_ACC}}} \\
\bottomrule
\end{tabular}
\end{table}

\textbf{Observation.} LoRA recovery adds \textbf{{{LORA\_GAIN}}\%} accuracy.

\subsection{Comparison to A3 / AASVD}

\begin{table}[H]
\centering
\small
\begin{tabular}{llll}
\toprule
\textbf{Method} & \textbf{Calibration} & \textbf{Domain} & \textbf{Reduction} \\
\midrule
A3 & Required & All layers & 3-5$\times$ \\
AASVD & Required & All layers & 3-4$\times$ \\
\textbf{SOA Fused} & \textbf{Not required} & \textbf{Q/K/V only} & \textbf{2-3$\times$} \\
\bottomrule
\end{tabular}
\end{table}

SOA Fused trades broader compression for \textbf{zero-shot} applicability.

\subsection{Failure Modes}

\subsubsection{Cases Where It Fails}

\begin{table}[H]
\centering
\small
\begin{tabular}{lll}
\toprule
\textbf{Case} & \textbf{Reason} & \textbf{Response} \\
\midrule
$k$ too small & Reconstruction error > threshold & Fallback \\
Wrong nested detection & Output projection double-applied & Fixed by \texttt{is\_nested\_encoder} \\
Qwen hybrid layers & Skipped by design & -- \\
\bottomrule
\end{tabular}
\end{table}

\subsubsection{RoBERTa Stability}

RoBERTa's loss curve is \textbf{less smooth} than DistilBERT's. Hypothesis: 
nested attention structure requires more recovery.

\subsection{Summary of Insights}

\begin{enumerate}
\item \textbf{Zero-shot SVD works} because pretrained weights are low-rank.
\item \textbf{$W_O$ matters} -- factorizing it would compound error.
\item \textbf{LoRA recovery adds {{LORA\_GAIN}}\%.}
\item \textbf{Rank 128 is a sweet spot} -- 3$\times$ reduction, minimal loss.
\item \textbf{No calibration is a major practical advantage.}

\end{enumerate}
